\section{Kræfter}

\[
    N = kg \cdot \frac{m}{s^2}
\]

\subsubsection{Newtons Love}
\begin{enumerate}
    \item Et objekt ændrer kun hastighed eller retning hvis det påvirkes af en ekstern kraft.
    \item $F = ma$ (med retning $\vec F = m\vec a$).
    \item For hver handling er der en lige og modsat reaktion.
\end{enumerate}

Den samlede eksterne kraft på et objekt er summen af all eksterne kræfter der påvirker det.

\subsubsection{Vægt}
Vægt defineres $\vec w = m \cdot \vec g$ i vektorform of $w = m \cdot g$ i skalarform.

\subsubsection{Normalkraft}
Normalkraften virker vinkelret på en overflade. 
\[
    N = mg \cos(\theta) \hspace{6em} \text{(På en horisontal overflade er } N = mg\text{)}
\]

\subsubsection{Snor-kraft}
Snor-kraften er kraften langs snoren. 
\[
    T = m \cdot a \sin(\theta) \hspace{6em} \text{(Ved 90° vil } T = mg\text{)}
\]

En vægt midt på en horisontal snor giver snor-kraften:
\[
    T = \frac{mg}{2 \sin(\theta)}
\]
Hvor $\theta$ er vinklen mellem snoren og den horisontale linje.

\subsubsection{Fjeder-kræft}
Fjeder-kræften modarbejder forlængelsen og kompression.
\[
    F = k \cdot x \kildetag{Hookes lov}
\]
Hvor $k$ er fjederkonstanten og $x$ er fjederens forlængelse eller kompression.

\subsection{Friktion}
Statisk friktion $f_s$ er friktionen før et objekt bevæger sig, kinetisk friktion $f_k$ er friktionen mens det bevæger sig.
\[
    f_s \le \mu_s N \hspace{6em} f_k = \mu_k N
\]
Hvor $\mu_s$ og $\mu_k$ er de statiske og kinetiske friktionskoefficienter, og $N$ er normalkraften.

\[
    \mu_k \le \mu_s
\]


\subsubsection{Friktion på hældninger}
Friktionskraften på en hældning er:
\[
    f = \mu_k \cdot m \cdot g \cdot \cos(\theta)
\]
Og accelerationen på en hældning er:
\[
    a = g \cdot (\sin(\theta) - \mu_k \cdot \cos(\theta))
\]
Hvis et objekt bevæger sig ned ad en hældning med konstant fart, er friktions-koeffienten:
\[
    \mu_k = \tan(\theta)
\]

\subsection{Centripetal kraft}
Centripetal kraften er den kraft der kræves for at holde et objekt i en cirkulær bane, og peger mod centrum af cirklen.
\[
    F_c = m \cdot a_c = m \cdot r \cdot \omega^2 = m \frac{v^2}{r}
\]

\subsection{Krængede kurver}
Bevægelse langs en krænget kurve giver en centripetal kraft der peger mod indersiden af kurven, uden at bruge friktion til at opnå det.

\subsubsection{Perfekt Krængning}

En perfekt krængning er en krænget kurve hvor ingen friktion er nødvendig for at holde objektet på banen. Sammenhængen mellem krængningsvinklen $\theta$ og hastigheden $v$ er:
\[
    \tan(\theta) = \frac{v^2}{r \cdot g}
\]

\subsection{Luftmodstand}
Luftmodstand virker mod bevægelsesretningen af objekter der bevæger sig gennem væsker.
\[
    F_D = \frac{1}{2} \cdot \rho \cdot v^2 \cdot C \cdot A
\]
Hvor $\rho =$ væskens densitet, $v =$ hastighed, $C =$ luftmodstands-koefficienten, og $A =$ tværsnitsarealet af objektet.

\subsubsection{Terminal-hastighed}
Terminal-hastigheden er hastigheden påkrævet for $F_D = mg$:
\[
    v_t = \sqrt{\frac{2mg}{\rho C A}}
\]

\subsubsection{Stoke's Lov}
Stoke's lov beskriver luftmodstand for små objekter i en væske:
\[
    F_D = 6 \pi \eta r v
\]
Hvor $\eta =$ væskens viskositet, $r =$ radius af objektet, og $v =$ hastighed.